\documentclass[runningheads]{llncs}

% ---------------------------------------------------------------
% ECCV 2026 package
% TODO REVIEW: Insert your submission number below by replacing '*****'
% TODO FINAL: Comment out the following line for the camera-ready version
\usepackage[review,year=2026,ID=*****]{eccv}
% TODO FINAL: Un-comment the following line for the camera-ready version
%\usepackage{eccv}

% ---------------------------------------------------------------
% Other packages
\usepackage{eccvabbrv}
\usepackage{graphicx}
\usepackage{booktabs}
\usepackage{amsmath}
\usepackage{amssymb}
\usepackage{multirow}
\usepackage{subcaption}
\usepackage[accsupp]{axessibility}

% ---------------------------------------------------------------
% Hyperref
% TODO FINAL: Comment out the following line for the camera-ready version
%\usepackage[pagebackref,breaklinks,colorlinks,citecolor=eccvblue]{hyperref}
% TODO FINAL: Un-comment the following line for the camera-ready version
\usepackage{hyperref}

\usepackage{orcidlink}

% Custom commands
\newcommand{\safe}{\textsc{Safe}}
\newcommand{\preffn}{Pre-FFN}
\newcommand{\kvaugment}{KV-Aug}

\begin{document}

% ---------------------------------------------------------------
\title{SAFE: Sensory Adapter Fusion for Enhanced\\Multimodal Understanding}

\titlerunning{SAFE: Sensory Adapter Fusion}

% TODO FINAL: Replace with your author list
\author{Anonymous ECCV 2026 Submission}
\authorrunning{Anonymous}
\institute{Anonymous Institution}

\maketitle

% ===============================================================
\begin{abstract}
Large vision-language models (LVLMs) excel at image-text reasoning but remain limited to their native modalities.
Extending these models to additional sensory inputs---audio, 3D point clouds, and beyond---typically requires costly retraining or modality-specific architectures.
We introduce \safe{} (\textbf{S}ensory \textbf{A}dapter \textbf{F}usion for \textbf{E}nhanced multimodal understanding), a lightweight, modality-agnostic framework that composes new sensory streams into frozen LVLMs via residual adapter fusion.
\safe{} projects arbitrary sensory inputs into the LVLM's hidden space through small learned adapters and fuses them as additive residuals at selected transformer layers, preserving the model's original vision-language capabilities while enabling new modalities with minimal trainable parameters.
We evaluate \safe{} on audio classification (ESC-50), 3D shape recognition (ModelNet40), and audio-visual question answering (MUSIC-AVQA), demonstrating that residual fusion achieves competitive or superior performance to modality-specific baselines while enabling true cross-modal composition in a single forward pass.

\keywords{Multimodal Learning \and Adapter Fusion \and Vision-Language Models \and Audio-Visual Understanding \and 3D Understanding}
\end{abstract}

% ===============================================================
\section{Introduction}
\label{sec:intro}

% TODO: Write introduction
% Key points:
% - LVLMs are powerful but locked to vision+language
% - Adding new modalities is expensive (full retraining) or brittle (separate models)
% - We propose SAFE: lightweight adapter fusion into frozen LVLMs
% - Key insight: additive residual injection at transformer layers
% - Works across modalities: audio, 3D, and their compositions with vision
% - Contributions: (1) SAFE framework, (2) Pre-FFN fusion mechanism, (3) composition results

% ===============================================================
\section{Related Work}
\label{sec:related}

% TODO: Write related work
% Sections to cover:
% 2.1 Large Vision-Language Models (LLaVA, InstructBLIP, etc.)
% 2.2 Multimodal Adapters and Fusion (LoRA, prefix tuning, adapter layers)
% 2.3 Audio-Visual Learning (LAVISH, APE, CAV-MAE)
% 2.4 3D Understanding with Language Models (PointBERT, Point-LLM, etc.)
% 2.5 Modality Composition

\subsection{Large Vision-Language Models}

\subsection{Multimodal Adapter Methods}

\subsection{Audio-Visual Learning}

\subsection{3D Understanding with Language Models}

% ===============================================================
\section{Method}
\label{sec:method}

% TODO: Write method section
% 3.1 Overview / Problem Formulation
% 3.2 Sensory Encoder and Projection
% 3.3 Pre-FFN Residual Fusion
% 3.4 Modality Composition
% 3.5 Training Objective

\subsection{Overview}
\label{sec:method-overview}

% Architecture overview figure should go here
% \begin{figure}[tb]
%   \centering
%   \includegraphics[width=\linewidth]{figures/safe_architecture}
%   \caption{Overview of \safe{}.}
%   \label{fig:architecture}
% \end{figure}

\subsection{Sensory Encoder and Projection}
\label{sec:method-encoder}

\subsection{\preffn{} Residual Fusion}
\label{sec:method-preffn}

\subsection{Modality Composition}
\label{sec:method-composition}

\subsection{Training Objective}
\label{sec:method-training}

% ===============================================================
\section{Experiments}
\label{sec:experiments}

\subsection{Experimental Setup}
\label{sec:exp-setup}

% TODO: Datasets, metrics, implementation details
% Datasets: ESC-50, ModelNet40, MUSIC-AVQA
% Base model: LLaVA-1.5-13B (frozen)
% Encoders: CLAP (audio), PointBERT (3D)

\subsection{Audio Classification: ESC-50}
\label{sec:exp-esc50}

% TODO: ESC-50 results
% 5-fold CV, comparison to SOTA
% Current result: 97.35% accuracy

\subsection{3D Shape Recognition: ModelNet40}
\label{sec:exp-modelnet}

% TODO: ModelNet40 results
% Classification accuracy, comparison to PointBERT/Point-LLM

\subsection{Audio-Visual Question Answering: MUSIC-AVQA}
\label{sec:exp-avqa}

% TODO: AVQA composition results
% This is the KEY experiment for the paper
% Conditions: both (composition), audio-only, image-only
% Metrics: Exact Match, Token F1, per-question-type breakdown

\subsection{Ablation Studies}
\label{sec:exp-ablations}

% TODO: Ablation studies
% - Fusion layer selection
% - Number of projected tokens
% - Fusion mechanism (Pre-FFN vs alternatives)
% - Encoder freezing strategy

% ===============================================================
\section{Analysis}
\label{sec:analysis}

% TODO: Deeper analysis
% - Visualization of attention patterns
% - What does composition learn?
% - Failure cases
% - Computational cost analysis

% ===============================================================
\section{Conclusion}
\label{sec:conclusion}

% TODO: Write conclusion

% ===============================================================
\section*{Acknowledgements}
% TODO FINAL: Add acknowledgements (omit for review)

% ---- Bibliography ----
\bibliographystyle{splncs04}
\bibliography{main}

\end{document}
